The previous chapter introduced a formal definition of a language and of a language model as a distribution over strings.
We now delve into the question of how we can learn a language model from finite data.

\paragraph{A note on terminology.}
Unfortunately, we will again encounter some ambiguous terminology.
%We defined a language model simply as some valid discrete probability distribution (or probability space) over a set of finite strings or sentences.
In \cref{sec:lm-tightness}, we explicitly define a language model as a valid probability distribution over $\kleene{\str}$, the Kleene closure of some alphabet $\alphabet$, which means that $\sum_{\str \in \kleene{\alphabet}} \pLM\left(\str\right) = 1$. 
As we will see later, this means that the model is \emph{tight}, whereas it is \emph{non-tight} if $\sum_{\str \in \kleene{\alphabet}} \pLM\left(\str\right) < 1$. 
% \ryan{reference previous section, we defined a language model as a valid probability distribution over the $\kleene{\alphabet}$, the Kleene closure of some alphabet $\alphabet$.}
Definitionally, then, all language models are tight.
However, it is standard in the literature to refer to many non-tight language models as language models well.
We pardon in advance the ambiguity that this introduces.
Over the course of the notes, we attempt to stick to the convention that the term ``language model'' without qualification only refers to a tight language model whereas a ``non-tight language model'' is used to refer to a language model in the more colloquial sense.
Linguistically, tight is acting as a non-intersective adjective.
Just as in English, where a fake gun is not a gun, so too in our course notes a non-tight language model is not a language model.
%That is, we will somewhat irresponsibly refer to any (parametrized) computational model of language as a language model as well.
This distinction does in fact matter.
On one hand, we can prove that many language models whose parameters are estimated from data, e.g., a finite-state language model estimated by means of maximum-likelihood estimation, are, in fact, tight.
On the other hand, we can show that this is \emph{not} true in general, i.e., \emph{not} language models estimated from data will be tight.
For instance, a recurrent neural network language model estimated through gradient descent may not be tight \citep{Chen2018}.
%(e.g., a recurrent neural network or an \ngram{} model) are \emph{not} strictly language models as we have defined them in \cref{def:language-model}.

\subsection{The distinguished symbols $\bos$ and $\eos$.}

We will continue with the assumption that we are modeling distribution over $\kleene{\alphabet}$, the Kleene closure of an alphabet $\alphabet$.
As we will see later, we will often define language models in an \emph{autoregressive} way, where the probability distribution over sentences will be broken into a sequence of conditional distributions over the next possible symbol given the previous symbols.
To be able to define the distribution over the first symbol in a sentence (given no context) as well as to be able to \emph{finish} a sentence without generating symbols in perpetuity, we augment the alphabet of symbols $\alphabet$ with two special symbols not in the original alphabet.
% However, to make the modeling more practical, we augment the alphabet $\alphabet$ with two special symbols.
\ryan{This is a poor way of putting it. It isn't about practicality. It's about the fact that we seek to estimate an autoregressive model, so it's necessary. \response{anej} Is this better?}

These are the \defn{\underline{b}eginning \underline{o}f \underline{s}entence} ($\bos$) and \defn{\underline{e}nd \underline{o}f \underline{s}entence} ($\eos$) symols.
For a string $\stry = \sym_1 \cdots \sym_\strlen$, we will suggestively denote $\sym_{0} \defeq \bos$ and $\sym_{\strlen + 1} \defeq \eos$ (even though $\bos, \eos \notin \alphabet$!).
They will prove useful when \emph{modeling} strings with computational models,\ryan{We should do a search-and-replace?} as we will see shortly.
Informally, you can think of the $\bos$ symbol as marking the beginning of the string, and the $\eos$ symbol as denoting the end of the string (or even terminating its generation, as we will see later).
Most of the time when working with computational models, we will augment $\alphabet$ with $\eos$: we define the augmented alphabet as $\eosalphabet \defeq \alphabet \cup \{\eos\}$ and assume that any computational model $\pM$ works over the augmented alphabet $\eosalphabet$. 

\subsection{Normalization of Language Models}

% A language model is therefore simply a probability distribution over $\kleene{\alphabet}$---it assigns any possible string its probability.
Given a language model $\pLM$, we can write for $\str = \sym_1 \ldots \sym_\strlen \in \kleene{\alphabet}$
\begin{equation}
    \pLM\left(\str\right) = \prod_{\idx = 1}^{\strlen}\pLM\left(\sym_\idx \mid \str_{<\idx}\right)
\end{equation}
using the \emph{chain rule} of probability.

When specifying $\pM$, we have two fundamental options.
Depending on whether we model $\pM\left(\str\right)$ for each string $\str \in \kleene{\eosalphabet}$ \emph{directly} or we model \emph{individual} conditional probabilities $\pM\left(\sym_\idx \mid \str_{<\idx}\right)$ we distinguish \emph{globally} and \emph{locally} normalized models.

\begin{definition}  \label{def:globally-normalized-model}
    In a \defn{globally normalized model} the probability $\pM\left(\str\right)$ is modeled as 
    \begin{equation} \label{eq:globally-normalized-model}
        \pM\left(\str\right) \defeq \softmax\left(\scoref\left(\str\right)\right) = \frac{\exp\left[\scoref\left(\str\right)\right]}{\sum_{\str' \in \kleene{\eosalphabet}}\exp\left[\scoref\left(\str'\right)\right]},
    \end{equation}
    where $\scoref$ is some real-valued scoring function that determines the probability of that string given by $\pM$ through the $\softmax$ function.
    The denominator $\normConstant_G \defeq \sum_{\str' \in \kleene{\eosalphabet}}\exp\left[\scoref\left(\str'\right)\right]$ is called the \defn{normalization constant}.
\end{definition}
Notice that the normalization constant is computed by summing over a countably infinite set of all possible strings. 
This means that the sum can in general \emph{diverge} and the quantity in \cref{eq:globally-normalized-model} may not be defined. 
This motivates the following definition.
\begin{definition}[Normalizable model]
    We say that a globally normalized model is \defn{normalizable} if the quantity $\normConstant_G$ in \cref{eq:globally-normalized-model} is finite.
\end{definition}

Net, we define locally-normalized models.
\begin{definition} \label{def:locally-normalized-model}
    In a \defn{locally normalized model} (also called an \defn{autoregressive model}) the probability $\pM\left(\str\right)$ is modeled using a sequence of \emph{conditional} distributions as
    \begin{align} \label{eq:locally-normalized-model}
        \pM\left(\str\right) \nonumber
        \defeq \prod_{\idx = 1}^{\strlen + 1}\pM\left(\sym_\idx\mid\str_{<\idx}\right) 
        &\defeq \prod_{\idx = 1}^{\strlen + 1} \softmax\left(\scoref\left(\sym_\idx \mid \str_{<\idx}\right)\right) \\ 
        &= \prod_{\idx = 1}^{\strlen + 1}\frac{\exp{\left[\scoref\left(\sym_\idx, \str_{<\idx}\right)\right]}}{\sum_{\sym_\idx' \in \eosalphabet}\exp{\left[\scoref\left(\sym_\idx', \str_{<\idx}\right)\right]}},
    \end{align}
    where we overload $\scoref$ to be some real-valued scoring function that defines the probability of the next symbol $\sym$ \emph{given} the context $\str_{<\idx}$ through the $\softmax$ function.
    Locally normalized models define multiple normalization constants $\normConstant_L^\idx \defeq \sum_{\sym_\idx' \in \eosalphabet}\exp{\left[\scoref\left(\sym_\idx', \str_{<\idx}\right)\right]}$.
\end{definition}

The names naturally come from the way the distributions in the two families are normalized: whereas globally normalized models are normalized by summing over the entire (infinite) space of strings, locally normalized models define a sequence of \emph{conditional distributions} and make use of the chain rule of probability to define the joint probability of a whole string.

\Anej{Add examples of a locally and globally normalized model (prefix tree-like in NLP slides}

As mentioned, if a model is globally normalized, its normalization constant might not be finite.
Even if it is finite, however, the fact that $\kleene{\eosalphabet}$ is \emph{infinite} means that we cannot generally compute $\normConstant_G$ in a \emph{tractable} way.
In fact, there are no general-purpose algorithms for this.
Because of that, locally normalized models which specify distributions over strings $\pdens\left(\sym_1 \ldots \sym_\strlen\right)$ in terms of their conditional probabilities $\pdens(\sym_\idx \mid \str_{<\idx})$ for $\idx \in \left[\strlen\right]$ have become the standard in NLP literature.
This reduces the problem of having to normalize the distribution over an infinite set $\kleene{\alphabet}$ to the problem of modeling the \emph{next possible symbol} $\sym_\idx$ given the symbols seen so far $\str_{<\idx}$.

All is not lost, however, for globally normalized models.
By making appropriate assumptions on the distribution $\pLM$, we can sometimes efficiently compute the normalizing constant $\normConstant_G$ as well. 
An example of such models is finite-state language models, which will introduce in \cref{sec:n-gram-models}.

Using the chain rule of probability, we can easily show the following theorem.
\begin{theorem} \label{thm:globally-normalized-is-locally-normalized}
Let $\pM$ be a globally-normalized language model.
Then there exists a locally-normalized language model $\pM'$ such that $\pM'\left(\str\right) = \pM\left(\str\right)$ for all $\str \in \kleene{\eosalphabet}$.
\end{theorem}
\begin{proof}
    We define the individual conditional probability distributions of the locally normalized-model $\pM'$ as:
    \begin{equation}
        \pM'\left(\sym\mid\str\right) = \frac{\pM'\left(\str\sym\right)}{\pM'\left(\str\right)}
    \end{equation}
    for $\sym \in \eosalphabet$ and $\str \in \kleene{\eosalphabet}$ wuch that $\pdens\left(\str\right) > 0$.
    Using axiom $2)$ from the definition of a probability measure (cf. \cref{def:probability-measure}), it is easy to see that the conditional probabilities defined this way indeed define a locally-normalized model.
\end{proof}

An advantage of globally normalized models is that they, by definition, always define a \emph{valid} probability space over $\eosalphabet$.
Although this might be counterintuitive at first, the same cannot be said for autoregressive models---in this sense, locally-normalized ``language models'' might not even be language models!
(This is the issue with the terminology we brought up earlier).
We explore this issue in much more detail in the next section.

\subsection{A General Framework} \label{sec:general-framework}
Most of the models we will encounter in the course will be \emph{locally normalized}.
In fact, many of them will define the conditional probability distribution over the next symbol, $\pM\left(\sym\mid\str\right)$ in a particular way.
It, therefore, makes sense to discuss this common feature upfront and then reflect on how it is implemented in particular models we will consider later.

We start by associating a real-valued column vector $\embedSym \in \R^\embeddim$ with each $\sym \in \eosalphabet$.
Additionally, we assume the existence of an \defn{encoding function} $\enc : \kleene{\alphabet} \rightarrow \Rd$ that maps every context $\str \in \kleene{\alphabet}$ to a real-valued vector $\enc\left(\str\right) \in \Rd$ and a real-valued scalar $\biasStr \in \R$. 
Then, we define the conditional probabilities $\pM\left(\sym\mid\str\right)$ as:
\begin{align} \label{eq:softmax}
\pM(\str) \defeq \prod_{\tstep=1}^{\strlen + 1} \pM(\sym_\tstep \mid \str_{<\tstep}) &= \prod_{\tstep=1}^{\strlen + 1} \softmax\left(\embedSymt^{\top} \enc(\str_{<\tstep}) + \biasStr\right) \\
&= \prod_{\tstep=1}^{\strlen + 1} \frac{\exp\left(\embedSymt^{\top} \enc(\str_{<\tstep}) + \biasStr \right) }{\sum_{\sym' \in \eosalphabet} \exp\left(\embedSymPrime^{\top} \enc(\str_{<\tstep}) + \biasStr\right)} 
\end{align}
where $\sym_0 \defeq \bos$ and $\sym_{\strlen + 1} \defeq \eos$.

Intuitively, we want $\embedSym$ and $\enc\left(\str\right)$ to capture the information needed to determine how probable individual symbols $\sym$ are given the context $\str$. 
This makes these embeddings crucial for accurate language modeling. 
Indeed, much of the discussion in class will center around how to build good embedding functions and embeddings of individual symbols. 
% \ryan{Mention that this is a softmax}
\ryan{Discuss the rule of exp}
\Ryan{Suggested Proof: Equivariance of the softmax?}
